\documentclass[10pt,a4paper]{amsart}
\usepackage{amsmath,amssymb,amsthm,graphicx}
% Retained minimal preamble for the Stacks Project chapter more-morphisms.tex (More on Morphisms).
% Verbatim excerpt of the chapter preamble: Xy-pic setup (native preamble lines 14-18) and the native
% theorem declarations with the mathematical macros (native preamble lines 154-217). Copied mechanically;
% nothing added, removed or redefined. The wrapper class and the amsmath/amssymb/amsthm/graphicx base
% packages are supplied by the assembler, so the class/externaldocument part of the native preamble is omitted.
% For commutative diagrams we use Xy-pic
\usepackage[all]{xy}

% We use 2cell for 2-commutative diagrams.
\xyoption{2cell}
% Theorem environments.
%
\theoremstyle{plain}
\newtheorem{theorem}[subsection]{Theorem}
\newtheorem{proposition}[subsection]{Proposition}
\newtheorem{lemma}[subsection]{Lemma}

\theoremstyle{definition}
\newtheorem{definition}[subsection]{Definition}
\newtheorem{example}[subsection]{Example}
\newtheorem{exercise}[subsection]{Exercise}
\newtheorem{situation}[subsection]{Situation}

\theoremstyle{remark}
\newtheorem{remark}[subsection]{Remark}
\newtheorem{remarks}[subsection]{Remarks}

\numberwithin{equation}{subsection}

% Macros
%
\def\lim{\mathop{\mathrm{lim}}\nolimits}
\def\colim{\mathop{\mathrm{colim}}\nolimits}
\def\Spec{\mathop{\mathrm{Spec}}}
\def\Hom{\mathop{\mathrm{Hom}}\nolimits}
\def\Ext{\mathop{\mathrm{Ext}}\nolimits}
\def\SheafHom{\mathop{\mathcal{H}\!\mathit{om}}\nolimits}
\def\SheafExt{\mathop{\mathcal{E}\!\mathit{xt}}\nolimits}
\def\Sch{\mathit{Sch}}
\def\Mor{\mathop{\mathrm{Mor}}\nolimits}
\def\Ob{\mathop{\mathrm{Ob}}\nolimits}
\def\Sh{\mathop{\mathit{Sh}}\nolimits}
\def\NL{\mathop{N\!L}\nolimits}
\def\CH{\mathop{\mathrm{CH}}\nolimits}
\def\proetale{{pro\text{-}\acute{e}tale}}
\def\etale{{\acute{e}tale}}
\def\QCoh{\mathit{QCoh}}
\def\Ker{\mathop{\mathrm{Ker}}}
\def\Im{\mathop{\mathrm{Im}}}
\def\Coker{\mathop{\mathrm{Coker}}}
\def\Coim{\mathop{\mathrm{Coim}}}

% Boxtimes
%
\DeclareMathSymbol{\boxtimes}{\mathbin}{AMSa}{"02}

%
% Macros for moduli stacks/spaces
%
\def\QCohstack{\mathcal{QC}\!\mathit{oh}}
\def\Cohstack{\mathcal{C}\!\mathit{oh}}
\def\Spacesstack{\mathcal{S}\!\mathit{paces}}
\def\Quotfunctor{\mathrm{Quot}}
\def\Hilbfunctor{\mathrm{Hilb}}
\def\Curvesstack{\mathcal{C}\!\mathit{urves}}
\def\Polarizedstack{\mathcal{P}\!\mathit{olarized}}
\def\Complexesstack{\mathcal{C}\!\mathit{omplexes}}
% \Pic is the operator that assigns to X its picard group, usage \Pic(X)
% \Picardstack_{X/B} denotes the Picard stack of X over B
% \Picardfunctor_{X/B} denotes the Picard functor of X over B
\def\Pic{\mathop{\mathrm{Pic}}\nolimits}
\def\Picardstack{\mathcal{P}\!\mathit{ic}}
\def\Picardfunctor{\mathrm{Pic}}
\def\Deformationcategory{\mathcal{D}\!\mathit{ef}}
\title{Finitely many Taylor coefficients of a formal local map are controlled by one finite type etale-like modification}
\author{The Stacks Project Authors}
\date{stacks-project:more-morphisms.tex:local-snapshot-20261011; GNU FDL 1.2 or later}
\begin{document}
\maketitle
\noindent\textit{Editorial note.} Statement span is the whole lemma environment lines 11529-11542; proof span is lines 11544-11568, the immediately following proof environment.\par\medskip
\noindent\textit{Editorial note.} The proof cites no label and no bibliography entry, so no external\_references are needed and the assembler adds no navigation stubs.\par\medskip
\noindent\textit{Editorial note.} The proof text is retained verbatim, including the original's use of the same symbol sigma for the map and for an arbitrary competitor map; no notational repair was applied.\par\medskip
\noindent\textit{Editorial note.} Difficulty H3 basis: the simultaneous control of all factorisations modulo m\_s\textasciicircum{}N is a genuine research-level finiteness construction, not a routine consequence of the definitions.\par\medskip
\noindent\textit{Editorial note.} extra\_packages is empty on purpose: the retained preamble already loads xy with options and a second declaration clashes.\par\medskip
\noindent\textit{Allowed standard prerequisites (editorial).}\begin{itemize}
\item Explicit prerequisite boundary. Noetherian local rings and their completions, finite type ring maps and their presentations by generators and relations, principal localisations, and the Taylor expansion of an element of the completion along a generating sequence of the maximal ideal. The control construction is retained in full.
\end{itemize}\medskip
\begin{lemma}
\label{lemma-control-agreement}
Let $T \to S$ be finite type morphisms of Noetherian schemes.
Let $t \in T$ map to $s \in S$ and let
$\sigma : \mathcal{O}_{T, t} \to \mathcal{O}_{S, s}^\wedge$
be a local $\mathcal{O}_{S, s}$-algebra map. For every $N \geq 1$
there exists a finite type morphism $(T', t') \to (T, t)$
such that $\sigma$ factors through
$\mathcal{O}_{T, t} \to \mathcal{O}_{T', t'}$
and such that for every local $\mathcal{O}_{S, s}$-algebra map
$\sigma' : \mathcal{O}_{T, t} \to \mathcal{O}_{S, s}^\wedge$
which factors through $\mathcal{O}_{T, t} \to \mathcal{O}_{T', t'}$
the maps $\sigma$ and $\sigma'$ agree modulo $\mathfrak m_s^N$.
\end{lemma}

\begin{proof}
We may assume $S$ and $T$ are affine. Say $S = \Spec(R)$ and
$T = \Spec(C)$. Let $c_1, \ldots, c_n \in C$ be generators
of $C$ as an $R$-algebra. Let $\mathfrak p \subset R$ be the
prime ideal corresponding to $s$. Say $\mathfrak p = (f_1, \ldots, f_m)$.
After replacing $R$ by a principal localization
(to clear denominators in $R_\mathfrak p$)
we may assume there exist
$r_1, \ldots, r_n \in R$ and $a_{i, I} \in \mathcal{O}_{S, s}^\wedge$
where $I = (i_1, \ldots, i_m)$ with $\sum i_j = N$ such that
$$
\sigma(c_i) = r_i + \sum\nolimits_I a_{i, I} f_1^{i_1} \ldots f_m^{i_m}
$$
in $\mathcal{O}_{S, s}^\wedge$. Then we consider
$$
C' = C[t_{i, I}]/
\left(c_i - r_i - \sum\nolimits_I t_{i, I} f_1^{i_1} \ldots f_m^{i_m}\right)
$$
with $\mathfrak p' = \mathfrak pC' + (t_{i, I})$ and factorization
of $\sigma : C \to \mathcal{O}_{S, s}^\wedge$ through $C'$
given by sending $t_{i, I}$ to $a_{i, I}$.
Taking $T' = \Spec(C')$ works because any $\sigma'$ as in the statement
of the lemma will send $c_i$ to $r_i$ modulo the maximal ideal
to the power $N$.
\end{proof}

\end{document}
